\subsection{The printed comparison question and the chronology gate}

\begin{result}{G.2.1}{The two printed passages requested in Lead H4}{Known; sources checked}
The Nice address contains, on printed page 435, the explicit reconstruction questions between a filtered Dieudonne module and the generic-fibre $\Q_p$-representation. The 1966 article contains the elliptic period-relation question in note 10 and the relative de Rham connection in note 13. These passages give exact printed antecedents for the respective questions in Parts A and B. They do not determine the date of the individual manuscript leaf 18:5.
\end{result}
\begin{proof}[Source comparison]
In \cite[p.~435, Remarks 2--3]{Grothendieck1970BT}, the base is a complete discrete valuation ring $A$ with perfect residue field $k$ of characteristic $p$ and fraction field $K$ of characteristic zero. The symbol $L$ denotes $\operatorname{Frac}W(k)$. The filtered object is an $L$-vector space $E$ with Frobenius and Verschiebung, together with $\Fil\subset E\otimes_LK$. Remark 2 asks: ``Comment peut-on expliciter $G$, et plus particulièrement sa fibre générique'' in terms of these data. The generic fibre is explicitly interpreted there as a finite-dimensional $\Q_p$-vector space with $\Gal(\overline K/K)$-action. Remark 3 asks how to ``reconstituer plus ou moins algébriquement le couple'' $(E,\Fil)$ from that representation. Thus the two arrows are exactly
\[
 (E,F,V,\Fil)\longleftrightarrow V_p(G_K),
\]
with the same scalar extension $L\to K$ required in B.1. The final paragraph of that page enlarges the question to longer filtrations and higher-dimensional $p$-adic bases. Part B gives the modern functors and their hypotheses. We read the scanned printed page, rather than inferring this content from the title of the address.

In \cite[note 10, p.~102]{Grothendieck1966}, the sentence beginning ``The only known general algebraic relation'' concerns the four elliptic periods and quasi-periods and displays Legendre's relation. The following sentences ask for algebraic independence and its extension to abelian varieties. The exact relation to A.4 is the period matrix of the two de Rham classes and two integral cycles; reversing the sign chosen for the second de Rham class reverses the quasi-period row and the written determinant convention. Note 13, explicitly dated July 1965 on p.~103, proposes the canonical integrable connection on relative de Rham cohomology. Its comparison is the complex local-system/Gauss--Manin comparison, whose relative $p$-adic counterpart is formulated in B.2. Neither passage specifies the reconstruction of a $p$-adic Galois representation from a filtered Dieudonne module.

The official Montpellier inventory \cite[cote 18 and note 16]{MontpellierInventory} dates the dossier approximately from 1969 to 1972, citing bibliography and reused Bourbaki versos from 1969--1972. The subsequently downloaded original scan confirms this cover date and the question on scan 5 \cite[cover and scan 5]{GrothendieckCote18Scan}. This is a date range for the dossier. Neither inspected page supplies a date for the particular leaf 18:5. The Nice talk took place in 1970 and its proceedings were printed in 1971. Therefore the printed text is established; the ordering of that text and the individual leaf remains an exact historical obligation. G.2.3 supplies a dated written antecedent from May 1966. Its relation to the individually undated leaf, and any claim of earliest priority, retain the exact chronology obligation stated there.
\end{proof}

\subsection{An arithmetic-specialization check in the elliptic search}

The search for the converse in 18:13 also led to the complete
18-page preprint \cite{BhowmikPathak2024}. We distinguish that
consulted version from its 23-page journal publication, whose
full text was not obtained. The following correction concerns
the statement printed in the consulted preprint only.

\begin{result}{G.2.2}{A specialization counterexample and the coefficient condition that repairs it}{Refuted; corrected statement proved}
The lower bound in Theorem~1.11 of the consulted version of
\cite{BhowmikPathak2024} is false under its printed hypotheses.
An exact counterexample is
\[
 \tau=i,\qquad f=E_4,\qquad g=j,\qquad
 \widetilde f=E_4+\frac{E_2E_6}{E_4}.
 \tag{G.2.2a}
\]

The following specialization statement holds. Let $f$ be a
nonzero meromorphic modular form of weight $k\ne0$, let $g$
be a nonconstant modular function, and write
\[
 \widetilde f=\sum_{r=0}^p f_r E_2^r,
 \qquad \operatorname{weight}(f_r)=w-2r.
\]
All forms and functions are for congruence subgroups and have
algebraic Fourier coefficients. At $\tau\in\mathbb H$, assume
that $f,g,f_0,\ldots,f_p$ are finite, $f(\tau)\ne0$, and
$f_r(\tau)\ne0$ for at least one $r\ge1$. Then, with
$q_\tau=\exp(2\pi i\tau)$,
\[
\begin{split}
 &\trdeg_{\Qb}\Qb
   \bigl(q_\tau,f(\tau),g(\tau),\widetilde f(\tau)\bigr)\\
 &\qquad=\trdeg_{\Qb}\Qb
   \bigl(q_\tau,E_2(\tau),E_4(\tau),E_6(\tau)\bigr)
   \ge3.
\end{split}
\]
Thus the nonvanishing condition is imposed on an actual
positive-depth coefficient at the point of evaluation.
\end{result}
\begin{proof}
The transformation law for $E_6$ at the fixed point of
$\tau\mapsto-1/\tau$ gives
$E_6(i)=i^6E_6(i)=-E_6(i)$, hence $E_6(i)=0$.
Since $E_4^3-E_6^2=1728\Delta$ and $\Delta$ has no zeros
in $\mathbb H$, $E_4(i)\ne0$ and $j(i)=1728$.
The expression in (G.2.2a) has weight four and depth one:
its coefficient of $E_2$ is the nonzero meromorphic
weight-two form $E_6/E_4$. Every Fourier coefficient is
rational. At infinity it begins $2-528q+O(q^2)$, so it
is nonconstant. At $i$ all values in question are finite,
and $f(i)=\widetilde f(i)=E_4(i)\ne0$. Consequently
\[
 \Qb\bigl(e^{-2\pi},f(i),g(i),\widetilde f(i)\bigr)
       =\Qb\bigl(e^{-2\pi},E_4(i)\bigr),
\]
whose transcendence degree is at most two. This verifies
every stated hypothesis of the consulted Theorem~1.11 and
contradicts its bound three. Its positive-depth coefficient
has become zero at the evaluation point. The rational
expansion and specialization identities are checked in
\texttt{check\_modular\_specialization.py}; the transcendence-degree
upper bound follows from the two displayed generators.

We prove the corrected statement. Pass to a common congruence
subgroup on which all the functions are defined. On its
compact modular curve, both $j$ and the nonconstant function
$g$ give finite morphisms to $\mathbf P^1$. The functions
and the curve are defined over $\Qb$; the function-field
description is \cite[Theorem 2.11, pp.10--11]{BhowmikPathak2024}.
At the point determined by $\tau$, finiteness implies that
the value of every regular modular function is algebraic
over $\Qb(g(\tau))$, and also over $\Qb(j(\tau))$.
This statement includes a transcendental value: its
scheme-theoretic image is then the generic point, and the
same finite extension of residue fields applies.

Put $\delta=\Delta(\tau)\ne0$. The modular function
$f^{12}/\Delta^k$ is finite and nonzero at $\tau$, so its
value is algebraic over $\Qb(g(\tau))$. The equality
\[
 f(\tau)^{12}
   =\bigl(f^{12}/\Delta^k\bigr)(\tau)\,\delta^k
\]
and $k\ne0$ make $\delta$ algebraic over
$\Qb(f(\tau),g(\tau))$. The same field therefore contains
$j(\tau)$, $E_4(\tau)$ and $E_6(\tau)$ algebraically, using
\[
 E_4(\tau)^3=j(\tau)\delta,\qquad
 E_6(\tau)^2=(j(\tau)-1728)\delta.
\]
Applying the same ratio construction to each $f_r$ shows
that its finite value is algebraic over this field; zero
values cause no problem. The specialized polynomial
\[
 \sum_{r=0}^p f_r(\tau)X^r-\widetilde f(\tau)
\]
is nonconstant by the stated coefficient condition, and
vanishes at $X=E_2(\tau)$. Hence $E_2(\tau)$ is algebraic
over $\Qb(f(\tau),g(\tau),\widetilde f(\tau))$.

Conversely the finite map given by $j$, followed by the same
twelfth-power ratio identities, makes $f(\tau),g(\tau)$
and all $f_r(\tau)$ algebraic over
$\Qb(E_4(\tau),E_6(\tau))$. The polynomial expression then
makes $\widetilde f(\tau)$ algebraic over the field obtained
by adjoining $E_2(\tau)$. These two algebraic extensions
prove the displayed equality after adjoining $q_\tau$.
Nesterenko's theorem \cite[Theorem 1]{EllNesterenko1996},
in the precise form used in A.4,
gives the lower bound three.
\end{proof}

This check does not replace $E_2$ by $E_2^*$ in that
theorem. The exact change of value is
$E_2^*(\tau)=E_2(\tau)-3/(\pi\operatorname{Im}\tau)$;
the additional term remains in the Hodge-splitting equation
of A.4.1. The failure of the attempted specialization
therefore leaves the non-CM algebraicity obligation in place.

Kumar's main theorem \cite[p.410, with Lemma 3 on p.412]{Kumar2021}
was also checked against that equation. Its premise is
$\Q$-linear dependence of the ordinary quasiperiods.
For the Hodge class the actual transformed row is
\[
 (\eta_1-h\omega_1,\eta_2-h\omega_2)
   =\frac{\pi}{\operatorname{Area}(\Lambda)}
        (\bar\omega_1,\bar\omega_2).
\]
Its ratio is $\bar\tau$, which is transcendental in the
non-CM algebraic-invariant case by Schneider. Thus this
transformed row satisfies neither the rational-coefficient
nor the algebraic-coefficient dependence premise. The
conclusion on Kumar's p.414 uses $\Qb$, whereas his main
theorem uses $\Q$; the original printed bars were visually
checked, since text extraction loses them.


\begin{result}{G.2.3}{The May 1966 written Tate-module question and its exact scope}{Primary source checked; individual leaf date retained}
A dated written formulation of the reconstruction question occurs in Grothendieck's May 1966 letter to Tate, before the Nice address. Its express starting case concerns a $p$-divisible group whose reduction is an extension of an ind-\'etale group by a toric $p$-divisible group; the final sentence then asks about more general groups. The letter supplies a written antecedent, while its catalogue status is unpublished. It does not determine the relative date of the individual undated leaf 18:5.
\end{result}
\begin{proof}[Source and comparison maps]
The first leaf of the original typescript carries the handwritten heading ``Thil, Mai 66''. The official inventory independently identifies the letter in cote~6 as dating from May 1966 \cite[cote 6]{MontpellierInventory}. On original numbered p.~30, \S3.2, the underlined question is to
\begin{quote}
``d\'eterminer, en termes d'une telle donn\'ee, quel est le module de Tate correspondant \`a la fibre g\'en\'erique de $\Phi$''.
\end{quote}
The preceding data are a Dieudonn\'e module over Witt vectors and a lift of its filtration over a complete DVR with perfect residue field; the reference to \S3.1 retains the stated reduction qualification \cite[\S\S3.1--3.2, pp.29--30]{AuditGrothendieckTate1966}. The source's first and last leaves were inspected visually. The letter is listed as unpublished in the correspondence catalogue \cite{AuditCSGCorrespondenceCatalogue}.

Here is the actual map from this degree-one question to the cohomological comparison of Part~B. Use its $K,K_0,\kappa$ and Tate conventions, take an abelian scheme $A/\mathcal O_K$ with ordinary reduction, and set $G=A[p^\infty]$. The canonical pairing gives an isomorphism of continuous $G_K$-modules
\[
 V:=H^1_{\mathrm{\acute et}}(A_{\overline K},\Q_p)
 \xrightarrow{\sim}(V_pG)^\vee.
\]
The contravariant cohomological Dieudonn\'e object is
\[
 D=H^1_{\mathrm{cris}}(A_\kappa/W(\kappa))[1/p],
\]
with its $\sigma$-semilinear Frobenius and the filtration on $D_K$
transported from $H^1_{\dR}(A_K/K)$ by crystalline/de Rham comparison.
The comparison in B.2 identifies $D_{\mathrm{cris}}(V)$ with $D$;
its inverse gives the actual continuous Galois representation
\[
 V_pG\simeq V_{\mathrm{cris}}(D)^\vee,
 \qquad
 V_{\mathrm{cris}}(D)=
 (D\otimes_{K_0}B_{\mathrm{cris}})^{\varphi=1}
 \cap\Fil^0(D_K\otimes_KB_{\dR}).
\]
Recovering the integral lattice $T_pG$ requires the integral data of Part~B as well. Thus this abelian ordinary case is a common instance of the letter and the later smooth-proper cohomological question, with the dualization explicitly retained.

For the covariant Dieudonn\'e convention, Cartier duality gives
\[
 D_{\mathrm{cov}}(G):=D_{\mathrm{contra}}(G^\vee)
 \simeq D_{\mathrm{contra}}(G)^\vee(-1),
 \qquad
 D_{\mathrm{cris}}(V_pG)=D_{\mathrm{cov}}(G)(1).
\]
On the underlying dual module the integral operators are
\[
 F_{\mathrm{cov}}(a)=\sigma\circ a\circ V_{\mathrm{contra}},
 \qquad
 V_{\mathrm{cov}}(a)=\sigma^{-1}\circ a\circ F_{\mathrm{contra}}.
\]
They are respectively $\sigma$- and $\sigma^{-1}$-semilinear and satisfy
$F_{\mathrm{cov}}V_{\mathrm{cov}}=V_{\mathrm{cov}}F_{\mathrm{cov}}=p$.
Rationally $F_{\mathrm{cov}}=p\varphi_{\mathrm{dual}}$, since
$K_0(-1)$ has Frobenius $p\sigma$.
The filtered duality is
\[
 \Fil^rD_{\mathrm{cov},K}
 =\operatorname{ann}(\Fil^{2-r}D_{\mathrm{contra},K}).
\]
This specifies the twist and annihilator filtration instead of assigning an unstated convention to the letter's symbol $M$.

Two other earlier printed sources have also been checked. The 1968 notes of Grothendieck's November--December 1966 lectures give the crystalline/de Rham comparison for a smooth proper Witt-vector lift \cite[\S7.2, pp.352--353]{AuditGrothendieckCrystals1968}. Tate's 1967 article asks for the higher-degree Hodge--Tate decomposition \cite[p.180, Remark]{Tate1967}. The latter receives the associated-graded realization of the filtered de Rham object in B.1.4; that explicitly identifies the information retained by its formulation. The inspected passages do not themselves specify the inverse filtered-Frobenius/Galois functor displayed above.

The particular question leaf is PDF scan~5, archive-numbered p.~4. Inspection at 180 dpi revealed no date or dated reused text on that leaf. The preceding scan is a tan title cover; the following handwritten sheet is smaller, and the subsequent scan is a reused SGA typescript. Their adjacency does not establish that the typescript is the verso of the question leaf. Consequently the dossier's approximate 1969--1972 range remains dossier-level evidence \cite[cover and scans 4--7]{GrothendieckCote18Scan}. The dated 1966 question is established, without asserting that it is the earliest formulation anywhere or deciding the individual leaf's chronological position.
\end{proof}


\subsection{What remains to be determined}

The results index records each numbered mathematical statement with its hypotheses, proof location, manuscript pages and checked inputs. The report records each requested item and each Lead separately. The following obligations remain part of the research, rather than conclusions supplied by the preceding results:
\begin{enumerate}[leftmargin=*]
\item For the non-CM elliptic curve $y^2=4x^3-4x-1$, determine whether its normalized $h$ is algebraic; then settle the same question for all non-CM curves over $\Qb$, including those without a real model. The exact formulas, the nonzero-value gate, the tensor-period reformulation and the tested analytic and arithmetic routes are in A.4.
\item Determine $\trdeg_\Q\Q(\tau,\bar\tau)$ for every elliptic curve over $\Qb$ not isogenous to its conjugate. A.3 identifies the remaining algebraic-relation shapes after excluding rational-function relations.
\item Determine the unconditional image-locus dimensions actually attained by all joint period points in A.2, beyond the cases in which joint period genericity has been established. The named period-conjecture implications remain implications only.
\item Characterize the complete essential image of geometric Hodge pairs in C.1. The abelian image and the explicit exclusions are proved; these do not enumerate all geometric Hodge structures.
\item Determine the general equality cases in D.1, and recover a precise field from the manuscript data in D.2 before assigning it a transcendence degree. The extra field introduced in Lead H5 is not a definition present in the supplied page.
\item Determine the plain integral torsor-class image and the full single-generator separation assertion of E.2 in the literal number-field scope. E.2.6 constructs the simultaneous integral comparison on the CM-abelian and Artin diagram for two places of a quartic field; extension to the entire smooth-projective diagram remains to be determined. E.2.7 gives the real-Frobenius obstruction to that extension as a comparison with the additional specified involution. The path construction over transcendence-degree-one fields does not transfer its refutation to number fields. The exact forgetful map from enriched structures to plain torsor classes is retained in E.3.
\item Finish the source and hypothesis checks identified individually by the remaining-claim inventory, including the chronological ordering in G.2.1. A missing original source is recorded as not obtained, even where an opened modern source proves the mathematical statement used.
\end{enumerate}

No general solution to these obligations is claimed. Every progress review records their state and the strongest proved gate, as required by the requested stopping rule.

